
\begin{exampletext}
  \ifcase\value{r6}\relax\or
    The war of lies was not a real war, but a simple legal error, compounded until records became a mess.
    \Glspl{scribe} called out for a distinction between fact and fable, which implied a consistent view of history.
    This meant that once a `fact' had been established, scholars changed all other records and accounts of history to become consistent.
    Every new `corrected' book had to make the long journey from city to city, where others would rewrite it, and pass copies along.

    Scholars guessed at what they did not know --- consistency eventually became more important than truth, and every spelling mistake and inkblot which made its way past the rushed editorial process dragged history in a new direction.

    The end of this era came with stories about how the initial confusion arose, which themselves had to create a consistent structure, so scholars began to explain the fake books entirely by reference to the adventurers who would pay incredible sums for books about long-lost treasures, without knowing if the book was true or false.
    As a result, early history contains more guesswork than real records.
  \or
    After the great elven crime, they began to cover up their lies with myriad false histories.
    Of course, writing simple falsehoods would fool no-one, so instead the elves produced thousands of copies of nearly-correct works, obscuring their own misdeeds in a mountain of conflicting histories, with names changing here, a king swapping his generals, or warriors slotting into the wrong century.

    For their final trick, the elves published books on how the confusion and lies all began with misguided rules for consistency, which began a fashion among scholars of laughing at the notion that any of their history has to contain `facts', or that it should `make sense when taken with the other stories'.
    The lie has worked, and to this day, nobody can say for sure exactly what the elves did that required so much work to cover up.
  \or
    Copying out complete books is a lot of work.
    A scribe must dip a pen into an inkwell around fifty-thousand times for a book of three hundred pages.
    If an inkwell takes around two hundred dips, then a book demands two hundred and fifty vials of ink.

    Many gnomes seek to optimize everything in their path, or anyone else's path.
    One of these gnomes, horrified with the work of writing, and too lazy to perform it, decided to create a complete solution to the problem, using alchemy.

    Life magic grew reeds for paper, and death magic helped to dry them once done.
    He created an apparition of a scribe with Light magic, then gave it the ability to touch and move objects with Force magic.
    More Life magic provided ink, and Force magic provided the inkwells (or a practical, if unnatural facsimile).

    Once his production-line stood complete, Mind magic gave the scribes the single wish of copying out books on history.
    They identified the books he owned which spoke about history, and began to copy them, day and night, seeking always to have \emph{more and more history}.

    These fake scribes, created through a web of alchemical theories, quickly found a hard limit to copying out history: they only had so much history to copy, so they could only create more copies, but could not perform their only desire --- to copy out \emph{more history}.

    The first day reached a fine conclusion as the `scribes' produced a copy of a pile of history books.
    On the second day, the gnome did not check the new books, so he had not noticed that they had \emph{additional} history.

    The study, practice, and essential meaning of history does not mean `truth'.
    Books on history do not need to be accurate.
    Books on history contain reasonable attempts at guessing our past.
    So in order to maximize the quantity of history, the `scribes' created more, with their best guesses.
    These scribes then created explanations for their own malformed views on history, and began to describe a period known as `the War of Lies`, where history became muddled due to a lot of sell-swords purchasing faked books on history.

    The books sold well, and for a short time, nobody noticed the inconsistencies between them.
    And by the time they did notice, children had been raised with false ideas, and nobody wanted to admit their families had learned a lot of nonsense.

    People learn much easier than they unlearn.

    One silver lining has come from this calamity --- the only thing the nonsensical history books kept constant --- a universal measurement of time, used by gnomish machines.
    Humans still use \glsentrylong{gmt} as the standard for measuring time to this day.
  \or
    Many a travelling adventurer crew made their money from researching fallen cities, then raiding locations which had the best loot.
    This produced a market for antiquated books which mentioned golden thrones, magical wonders, lochs of eternal strength, and trapped scrolls which reveal the answer to any question posed.

    Unethical scholars soon began creating a horde of `hidden histories' books, which looked authentic and contained subtle clues to glorious treasures.
    Of course, claiming the `lost city of gold' was over the hill meant angry adventurers returning soon after.
    The smartest of the con-artists always hinted at treasures lying in the most dangerous areas of the deep forest.

    Scholars `deduced' any historical details they did not know about, but since adventurers never read much about history, the books' accuracy hardly mattered.
    Sensationalism trumps truth when it comes to sales, so the writers relied on creating tales to cater to their readers' prejudices and hopes.
    And when angry adventurers returned from some desolate swamp, empty-handed, the charlatans could always blame their book's errors on `an elvish conspiracy'.
  \else
    Having one word for `story', or `history', or `account' never confused anyone talking about anything.
    But context-heavy languages work less well in writing.
    Having one word for `true', `simple', `yes', `good', and `fun' added another complication, as people simultaneously assessed books for their historical accuracy, and comedic timing.
    Critics would respond that books created `confusion', and that word can also mean `bad', or `false'.

    Many came to the same beautiful conclusion; they would simply use a written language designed for accuracy in writing, then translate all the books into the accurate language.
    Some used the `official' Dwarvish writing, others used the local Elvish, and a few made a bespoke language for writing.

    The choice of language implied deep political affiliations with a local noble, a god, or one of the new Philosophers (who came into being retroactively, through the process of translating books in a way that made sense).
    The communication difficulties became so ingrained that trade slowed along linguistic lines, and anyone who wrote things a lot felt the world getting smaller with each generation.
    The final spark came in Whiteplains' central market, when a trader wanted to repay her debt in feet.
    The fist-fight turned into civil war.
    Any noble capable of communicating with another banded together out of fear, and war began organically, since everyone figured it would and nobody could, officially, disagree.

    The \gls{tradeTongue} was still available of course, but reliance on it quickly suggested nefarious activities.
    It suggested an attempt to make strange and unintelligible enemies intelligible again.
    Nobles on all sides began to burn their original books as a public display of cutting ties with everyone who wrote things down wrong.
    Public displays, of course, mean people can speak, which let them decry anyone who wrote things which were `wrong', `confusing', `bad', or `complicated', in a single word, without worrying about particulars.

    Half the towns in \gls{fenestra} have a story about a cunning warrior-scholar from their area, who solved the communications mess and stopped the war.
    So that famous ending to the war of lies is itself part of the war of lies, but at least it was the end.

    Those days lie behind us now.
    This is the very start of history.
  \fi

\end{exampletext}

